\documentclass[aspectratio=169,10pt]{beamer}
\makeatletter\def\input@path{{../../estilo/}{../estilo/}{estilo/}}\makeatother
\usepackage{devmobile}
\modulo[Flutter]{01}{Ambiente de Desenvolvimento}{Flutter SDK, editores, emuladores e celular real, estrutura do projeto, pubspec, hot reload, DevTools e pacotes}

\begin{document}
\capa

\begin{frame}{Objetivos do módulo}
  \begin{conceito}[Ao final deste módulo você será capaz de]
    \begin{itemize}
      \item Explicar o que é o Flutter e como ele chega ao Android, ao iOS, à web e ao desktop.
      \item Instalar o Flutter SDK e diagnosticar o ambiente com \cod{flutter doctor}.
      \item Executar apps no emulador, no simulador iOS e no celular real.
      \item Entender a estrutura de um projeto Flutter e o arquivo \arq{pubspec.yaml}.
      \item Usar \emph{hot reload}, os modos de execução e o Flutter DevTools.
      \item Encontrar, avaliar e adicionar pacotes do \textbf{pub.dev}.
    \end{itemize}
  \end{conceito}
\end{frame}

\sumario

% =============================================================================
%  Fundamentos compartilhados --- Módulo 01: ferramentas de desenvolvimento
%  Incluído pelas trilhas Kotlin e Flutter com % =============================================================================
%  Fundamentos compartilhados --- Módulo 01: ferramentas de desenvolvimento
%  Incluído pelas trilhas Kotlin e Flutter com % =============================================================================
%  Fundamentos compartilhados --- Módulo 01: ferramentas de desenvolvimento
%  Incluído pelas trilhas Kotlin e Flutter com \input{../fundamentos/f01-ambiente}
% =============================================================================
\section{Fundamentos: ferramentas de desenvolvimento}

\begin{frame}{Do código-fonte ao programa}
  \centering
  \begin{tikzpicture}[node distance=5mm,every node/.style={font=\scriptsize,align=center,minimum height=10mm}]
    \node[caixa,text width=1.9cm] (src) {Código-fonte\\ \tiny texto legível};
    \node[caixapetroleo,right=of src,text width=1.9cm] (comp) {Compilador\\ \tiny verifica e traduz};
    \node[caixaambar,right=of comp,text width=2.2cm] (bc) {\emph{Bytecode} ou\\ código de máquina};
    \node[caixapetroleo,right=of bc,text width=1.9cm] (pk) {Empacotador\\ \tiny junta recursos};
    \node[caixaescura,right=of pk,text width=1.9cm] (app) {App\\ instalável};
    \foreach \a/\b in {src/comp,comp/bc,bc/pk,pk/app} \draw[seta] (\a) -- (\b);
  \end{tikzpicture}
  \vspace{4pt}
  \begin{columns}[T]
    \begin{column}{0.48\textwidth}
      \begin{definicao}[Compilador e interpretador]
        O \textbf{compilador} traduz o programa inteiro antes da execução e aponta erros cedo. O \textbf{interpretador} lê e executa o código aos poucos, durante a execução.
      \end{definicao}
      \begin{definicao}[Máquina virtual]
        Programa que executa um \emph{bytecode} intermediário em qualquer processador: a JVM e o ART (Kotlin) e a VM do Dart.
      \end{definicao}
    \end{column}
    \begin{column}{0.48\textwidth}
      \begin{definicao}[JIT e AOT]
        \textbf{JIT}: compila enquanto o programa roda; permite trocar código em execução (\emph{hot reload}).
        \textbf{AOT}: compila tudo antes de distribuir; inicia rápido e roda com desempenho previsível.
      \end{definicao}
      \begin{dica}
        As duas trilhas usam os dois: JIT para desenvolver com agilidade, AOT para o app que chega ao usuário (Módulo 15).
      \end{dica}
    \end{column}
  \end{columns}
\end{frame}

\begin{frame}{SDK, IDE e cadeia de ferramentas}
  \begin{columns}[T]
    \begin{column}{0.48\textwidth}
      \begin{definicao}[SDK --- \emph{Software Development Kit}]
        Pacote com bibliotecas, compiladores, ferramentas de linha de comando, documentação e imagens de emulador para desenvolver para uma plataforma.
      \end{definicao}
      \begin{definicao}[IDE --- ambiente integrado]
        Editor que reúne escrita de código, compilação, execução, depuração, análise e integração com Git.
      \end{definicao}
      \begin{analogia}
        O SDK é a caixa de ferramentas; a IDE é a bancada organizada onde você trabalha com elas.
      \end{analogia}
    \end{column}
    \begin{column}{0.5\textwidth}
      \scriptsize
      \begin{tabularx}{\linewidth}{@{}l L L@{}}
        \toprule
        & \cabecalho{Trilha Kotlin} & \cabecalho{Trilha Flutter} \\
        \midrule
        Linguagem & Kotlin (sobre o JDK) & Dart (Dart SDK) \\
        SDK da plataforma & Android SDK & Flutter SDK + Android SDK + Xcode \\
        IDE & Android Studio & VS Code ou Android Studio \\
        \emph{Build} & Gradle & \cod{flutter} CLI (usa Gradle e Xcode) \\
        Aparelho & ADB & ADB e \cod{flutter devices} \\
        Pacotes & Maven Central, Google Maven & pub.dev \\
        \bottomrule
      \end{tabularx}
    \end{column}
  \end{columns}
\end{frame}

\begin{frame}{Sistemas de \emph{build} e dependências}
  \begin{columns}[c]
    \begin{column}{0.46\textwidth}
      \centering
      \begin{tikzpicture}[node distance=4mm,every node/.style={font=\tiny,minimum width=1.6cm}]
        \node[caixaescura] (apk) {gerar pacote};
        \node[caixapetroleo,below left=5mm and -2mm of apk] (comp) {compilar};
        \node[caixapetroleo,below right=5mm and -2mm of apk] (res) {processar recursos};
        \node[caixa,below=5mm of comp] (dep) {baixar dependências};
        \node[caixa,below=5mm of res] (gen) {gerar código};
        \draw[seta] (comp) -- (apk); \draw[seta] (res) -- (apk);
        \draw[seta] (dep) -- (comp); \draw[seta] (gen) -- (comp);
      \end{tikzpicture}\\[2pt]
      {\tiny Grafo de tarefas: cada uma só roda depois das que ela usa.}
    \end{column}
    \begin{column}{0.52\textwidth}
      \begin{definicao}[Sistema de \emph{build}]
        Ferramenta que transforma o projeto em app executando \textbf{tarefas} organizadas num grafo de dependências. Só refaz o que mudou (\emph{build} incremental) e guarda resultados em cache.
      \end{definicao}
      \begin{definicao}[Dependência transitiva]
        A biblioteca que você usa depende de outras, que dependem de outras. O gerenciador resolve a árvore inteira e escolhe versões compatíveis.
      \end{definicao}
      \begin{atencao}
        Fixe as versões com arquivos de trava (catálogo de versões, \cod{pubspec.lock}) para que todos na equipe compilem o mesmo app.
      \end{atencao}
    \end{column}
  \end{columns}
\end{frame}

\begin{frame}{Emuladores, simuladores e aparelhos reais}
  \begin{columns}[T]
    \begin{column}{0.5\textwidth}
      \scriptsize
      \begin{tabularx}{\linewidth}{@{}l L L@{}}
        \toprule
        & \cabecalho{Vantagem} & \cabecalho{Limite} \\
        \midrule
        Emulador Android & vários tamanhos e versões, GPS e sensores simulados & sem Bluetooth, NFC e USB reais \\
        Simulador iOS & rápido, integrado ao Xcode & roda código do Mac; não imita o hardware \\
        Aparelho real & desempenho e hardware verdadeiros & cada modelo exige um aparelho \\
        \bottomrule
      \end{tabularx}
    \end{column}
    \begin{column}{0.48\textwidth}
      \begin{definicao}[Emulador e simulador]
        O \textbf{emulador} imita o aparelho, inclusive o processador e o sistema completo. O \textbf{simulador} apenas reproduz o comportamento do sistema, executando o app compilado para o computador.
      \end{definicao}
      \begin{pordentro}[Por que o emulador exige virtualização]
        O emulador Android roda uma imagem completa do sistema numa máquina virtual. Com aceleração por hardware (Intel VT-x, AMD-V, KVM no Linux, Hypervisor no macOS), as instruções rodam direto no processador e o emulador fica rápido.
      \end{pordentro}
    \end{column}
  \end{columns}
\end{frame}

\begin{frame}{Depuração}
  \begin{columns}[T]
    \begin{column}{0.5\textwidth}
      \begin{definicao}[Depurador]
        Ferramenta que pausa o programa num ponto escolhido (\emph{breakpoint}) para inspecionar variáveis, a pilha de chamadas e executar passo a passo.
      \end{definicao}
      \scriptsize
      \begin{tabularx}{\linewidth}{@{}l L@{}}
        \toprule
        \cabecalho{Comando} & \cabecalho{O que faz} \\
        \midrule
        \emph{Step over} & executa a linha e para na próxima \\
        \emph{Step into} & entra na função chamada \\
        \emph{Step out} & termina a função atual e volta \\
        \emph{Resume} & continua até o próximo \emph{breakpoint} \\
        \bottomrule
      \end{tabularx}
    \end{column}
    \begin{column}{0.48\textwidth}
      \begin{conceito}[Método científico da depuração]
        \begin{enumerate}
          \item Reproduza o problema de forma confiável.
          \item Formule uma hipótese sobre a causa.
          \item Teste a hipótese (\emph{breakpoint}, \emph{log}, teste automatizado).
          \item Corrija e escreva um teste que impeça a volta do defeito.
        \end{enumerate}
      \end{conceito}
      \begin{dica}
        Leia a mensagem de erro e a pilha de chamadas inteiras: a primeira linha do \textbf{seu} código costuma apontar a causa.
      \end{dica}
    \end{column}
  \end{columns}
\end{frame}

% =============================================================================
\section{Fundamentos: ferramentas de desenvolvimento}

\begin{frame}{Do código-fonte ao programa}
  \centering
  \begin{tikzpicture}[node distance=5mm,every node/.style={font=\scriptsize,align=center,minimum height=10mm}]
    \node[caixa,text width=1.9cm] (src) {Código-fonte\\ \tiny texto legível};
    \node[caixapetroleo,right=of src,text width=1.9cm] (comp) {Compilador\\ \tiny verifica e traduz};
    \node[caixaambar,right=of comp,text width=2.2cm] (bc) {\emph{Bytecode} ou\\ código de máquina};
    \node[caixapetroleo,right=of bc,text width=1.9cm] (pk) {Empacotador\\ \tiny junta recursos};
    \node[caixaescura,right=of pk,text width=1.9cm] (app) {App\\ instalável};
    \foreach \a/\b in {src/comp,comp/bc,bc/pk,pk/app} \draw[seta] (\a) -- (\b);
  \end{tikzpicture}
  \vspace{4pt}
  \begin{columns}[T]
    \begin{column}{0.48\textwidth}
      \begin{definicao}[Compilador e interpretador]
        O \textbf{compilador} traduz o programa inteiro antes da execução e aponta erros cedo. O \textbf{interpretador} lê e executa o código aos poucos, durante a execução.
      \end{definicao}
      \begin{definicao}[Máquina virtual]
        Programa que executa um \emph{bytecode} intermediário em qualquer processador: a JVM e o ART (Kotlin) e a VM do Dart.
      \end{definicao}
    \end{column}
    \begin{column}{0.48\textwidth}
      \begin{definicao}[JIT e AOT]
        \textbf{JIT}: compila enquanto o programa roda; permite trocar código em execução (\emph{hot reload}).
        \textbf{AOT}: compila tudo antes de distribuir; inicia rápido e roda com desempenho previsível.
      \end{definicao}
      \begin{dica}
        As duas trilhas usam os dois: JIT para desenvolver com agilidade, AOT para o app que chega ao usuário (Módulo 15).
      \end{dica}
    \end{column}
  \end{columns}
\end{frame}

\begin{frame}{SDK, IDE e cadeia de ferramentas}
  \begin{columns}[T]
    \begin{column}{0.48\textwidth}
      \begin{definicao}[SDK --- \emph{Software Development Kit}]
        Pacote com bibliotecas, compiladores, ferramentas de linha de comando, documentação e imagens de emulador para desenvolver para uma plataforma.
      \end{definicao}
      \begin{definicao}[IDE --- ambiente integrado]
        Editor que reúne escrita de código, compilação, execução, depuração, análise e integração com Git.
      \end{definicao}
      \begin{analogia}
        O SDK é a caixa de ferramentas; a IDE é a bancada organizada onde você trabalha com elas.
      \end{analogia}
    \end{column}
    \begin{column}{0.5\textwidth}
      \scriptsize
      \begin{tabularx}{\linewidth}{@{}l L L@{}}
        \toprule
        & \cabecalho{Trilha Kotlin} & \cabecalho{Trilha Flutter} \\
        \midrule
        Linguagem & Kotlin (sobre o JDK) & Dart (Dart SDK) \\
        SDK da plataforma & Android SDK & Flutter SDK + Android SDK + Xcode \\
        IDE & Android Studio & VS Code ou Android Studio \\
        \emph{Build} & Gradle & \cod{flutter} CLI (usa Gradle e Xcode) \\
        Aparelho & ADB & ADB e \cod{flutter devices} \\
        Pacotes & Maven Central, Google Maven & pub.dev \\
        \bottomrule
      \end{tabularx}
    \end{column}
  \end{columns}
\end{frame}

\begin{frame}{Sistemas de \emph{build} e dependências}
  \begin{columns}[c]
    \begin{column}{0.46\textwidth}
      \centering
      \begin{tikzpicture}[node distance=4mm,every node/.style={font=\tiny,minimum width=1.6cm}]
        \node[caixaescura] (apk) {gerar pacote};
        \node[caixapetroleo,below left=5mm and -2mm of apk] (comp) {compilar};
        \node[caixapetroleo,below right=5mm and -2mm of apk] (res) {processar recursos};
        \node[caixa,below=5mm of comp] (dep) {baixar dependências};
        \node[caixa,below=5mm of res] (gen) {gerar código};
        \draw[seta] (comp) -- (apk); \draw[seta] (res) -- (apk);
        \draw[seta] (dep) -- (comp); \draw[seta] (gen) -- (comp);
      \end{tikzpicture}\\[2pt]
      {\tiny Grafo de tarefas: cada uma só roda depois das que ela usa.}
    \end{column}
    \begin{column}{0.52\textwidth}
      \begin{definicao}[Sistema de \emph{build}]
        Ferramenta que transforma o projeto em app executando \textbf{tarefas} organizadas num grafo de dependências. Só refaz o que mudou (\emph{build} incremental) e guarda resultados em cache.
      \end{definicao}
      \begin{definicao}[Dependência transitiva]
        A biblioteca que você usa depende de outras, que dependem de outras. O gerenciador resolve a árvore inteira e escolhe versões compatíveis.
      \end{definicao}
      \begin{atencao}
        Fixe as versões com arquivos de trava (catálogo de versões, \cod{pubspec.lock}) para que todos na equipe compilem o mesmo app.
      \end{atencao}
    \end{column}
  \end{columns}
\end{frame}

\begin{frame}{Emuladores, simuladores e aparelhos reais}
  \begin{columns}[T]
    \begin{column}{0.5\textwidth}
      \scriptsize
      \begin{tabularx}{\linewidth}{@{}l L L@{}}
        \toprule
        & \cabecalho{Vantagem} & \cabecalho{Limite} \\
        \midrule
        Emulador Android & vários tamanhos e versões, GPS e sensores simulados & sem Bluetooth, NFC e USB reais \\
        Simulador iOS & rápido, integrado ao Xcode & roda código do Mac; não imita o hardware \\
        Aparelho real & desempenho e hardware verdadeiros & cada modelo exige um aparelho \\
        \bottomrule
      \end{tabularx}
    \end{column}
    \begin{column}{0.48\textwidth}
      \begin{definicao}[Emulador e simulador]
        O \textbf{emulador} imita o aparelho, inclusive o processador e o sistema completo. O \textbf{simulador} apenas reproduz o comportamento do sistema, executando o app compilado para o computador.
      \end{definicao}
      \begin{pordentro}[Por que o emulador exige virtualização]
        O emulador Android roda uma imagem completa do sistema numa máquina virtual. Com aceleração por hardware (Intel VT-x, AMD-V, KVM no Linux, Hypervisor no macOS), as instruções rodam direto no processador e o emulador fica rápido.
      \end{pordentro}
    \end{column}
  \end{columns}
\end{frame}

\begin{frame}{Depuração}
  \begin{columns}[T]
    \begin{column}{0.5\textwidth}
      \begin{definicao}[Depurador]
        Ferramenta que pausa o programa num ponto escolhido (\emph{breakpoint}) para inspecionar variáveis, a pilha de chamadas e executar passo a passo.
      \end{definicao}
      \scriptsize
      \begin{tabularx}{\linewidth}{@{}l L@{}}
        \toprule
        \cabecalho{Comando} & \cabecalho{O que faz} \\
        \midrule
        \emph{Step over} & executa a linha e para na próxima \\
        \emph{Step into} & entra na função chamada \\
        \emph{Step out} & termina a função atual e volta \\
        \emph{Resume} & continua até o próximo \emph{breakpoint} \\
        \bottomrule
      \end{tabularx}
    \end{column}
    \begin{column}{0.48\textwidth}
      \begin{conceito}[Método científico da depuração]
        \begin{enumerate}
          \item Reproduza o problema de forma confiável.
          \item Formule uma hipótese sobre a causa.
          \item Teste a hipótese (\emph{breakpoint}, \emph{log}, teste automatizado).
          \item Corrija e escreva um teste que impeça a volta do defeito.
        \end{enumerate}
      \end{conceito}
      \begin{dica}
        Leia a mensagem de erro e a pilha de chamadas inteiras: a primeira linha do \textbf{seu} código costuma apontar a causa.
      \end{dica}
    \end{column}
  \end{columns}
\end{frame}

% =============================================================================
\section{Fundamentos: ferramentas de desenvolvimento}

\begin{frame}{Do código-fonte ao programa}
  \centering
  \begin{tikzpicture}[node distance=5mm,every node/.style={font=\scriptsize,align=center,minimum height=10mm}]
    \node[caixa,text width=1.9cm] (src) {Código-fonte\\ \tiny texto legível};
    \node[caixapetroleo,right=of src,text width=1.9cm] (comp) {Compilador\\ \tiny verifica e traduz};
    \node[caixaambar,right=of comp,text width=2.2cm] (bc) {\emph{Bytecode} ou\\ código de máquina};
    \node[caixapetroleo,right=of bc,text width=1.9cm] (pk) {Empacotador\\ \tiny junta recursos};
    \node[caixaescura,right=of pk,text width=1.9cm] (app) {App\\ instalável};
    \foreach \a/\b in {src/comp,comp/bc,bc/pk,pk/app} \draw[seta] (\a) -- (\b);
  \end{tikzpicture}
  \vspace{4pt}
  \begin{columns}[T]
    \begin{column}{0.48\textwidth}
      \begin{definicao}[Compilador e interpretador]
        O \textbf{compilador} traduz o programa inteiro antes da execução e aponta erros cedo. O \textbf{interpretador} lê e executa o código aos poucos, durante a execução.
      \end{definicao}
      \begin{definicao}[Máquina virtual]
        Programa que executa um \emph{bytecode} intermediário em qualquer processador: a JVM e o ART (Kotlin) e a VM do Dart.
      \end{definicao}
    \end{column}
    \begin{column}{0.48\textwidth}
      \begin{definicao}[JIT e AOT]
        \textbf{JIT}: compila enquanto o programa roda; permite trocar código em execução (\emph{hot reload}).
        \textbf{AOT}: compila tudo antes de distribuir; inicia rápido e roda com desempenho previsível.
      \end{definicao}
      \begin{dica}
        As duas trilhas usam os dois: JIT para desenvolver com agilidade, AOT para o app que chega ao usuário (Módulo 15).
      \end{dica}
    \end{column}
  \end{columns}
\end{frame}

\begin{frame}{SDK, IDE e cadeia de ferramentas}
  \begin{columns}[T]
    \begin{column}{0.48\textwidth}
      \begin{definicao}[SDK --- \emph{Software Development Kit}]
        Pacote com bibliotecas, compiladores, ferramentas de linha de comando, documentação e imagens de emulador para desenvolver para uma plataforma.
      \end{definicao}
      \begin{definicao}[IDE --- ambiente integrado]
        Editor que reúne escrita de código, compilação, execução, depuração, análise e integração com Git.
      \end{definicao}
      \begin{analogia}
        O SDK é a caixa de ferramentas; a IDE é a bancada organizada onde você trabalha com elas.
      \end{analogia}
    \end{column}
    \begin{column}{0.5\textwidth}
      \scriptsize
      \begin{tabularx}{\linewidth}{@{}l L L@{}}
        \toprule
        & \cabecalho{Trilha Kotlin} & \cabecalho{Trilha Flutter} \\
        \midrule
        Linguagem & Kotlin (sobre o JDK) & Dart (Dart SDK) \\
        SDK da plataforma & Android SDK & Flutter SDK + Android SDK + Xcode \\
        IDE & Android Studio & VS Code ou Android Studio \\
        \emph{Build} & Gradle & \cod{flutter} CLI (usa Gradle e Xcode) \\
        Aparelho & ADB & ADB e \cod{flutter devices} \\
        Pacotes & Maven Central, Google Maven & pub.dev \\
        \bottomrule
      \end{tabularx}
    \end{column}
  \end{columns}
\end{frame}

\begin{frame}{Sistemas de \emph{build} e dependências}
  \begin{columns}[c]
    \begin{column}{0.46\textwidth}
      \centering
      \begin{tikzpicture}[node distance=4mm,every node/.style={font=\tiny,minimum width=1.6cm}]
        \node[caixaescura] (apk) {gerar pacote};
        \node[caixapetroleo,below left=5mm and -2mm of apk] (comp) {compilar};
        \node[caixapetroleo,below right=5mm and -2mm of apk] (res) {processar recursos};
        \node[caixa,below=5mm of comp] (dep) {baixar dependências};
        \node[caixa,below=5mm of res] (gen) {gerar código};
        \draw[seta] (comp) -- (apk); \draw[seta] (res) -- (apk);
        \draw[seta] (dep) -- (comp); \draw[seta] (gen) -- (comp);
      \end{tikzpicture}\\[2pt]
      {\tiny Grafo de tarefas: cada uma só roda depois das que ela usa.}
    \end{column}
    \begin{column}{0.52\textwidth}
      \begin{definicao}[Sistema de \emph{build}]
        Ferramenta que transforma o projeto em app executando \textbf{tarefas} organizadas num grafo de dependências. Só refaz o que mudou (\emph{build} incremental) e guarda resultados em cache.
      \end{definicao}
      \begin{definicao}[Dependência transitiva]
        A biblioteca que você usa depende de outras, que dependem de outras. O gerenciador resolve a árvore inteira e escolhe versões compatíveis.
      \end{definicao}
      \begin{atencao}
        Fixe as versões com arquivos de trava (catálogo de versões, \cod{pubspec.lock}) para que todos na equipe compilem o mesmo app.
      \end{atencao}
    \end{column}
  \end{columns}
\end{frame}

\begin{frame}{Emuladores, simuladores e aparelhos reais}
  \begin{columns}[T]
    \begin{column}{0.5\textwidth}
      \scriptsize
      \begin{tabularx}{\linewidth}{@{}l L L@{}}
        \toprule
        & \cabecalho{Vantagem} & \cabecalho{Limite} \\
        \midrule
        Emulador Android & vários tamanhos e versões, GPS e sensores simulados & sem Bluetooth, NFC e USB reais \\
        Simulador iOS & rápido, integrado ao Xcode & roda código do Mac; não imita o hardware \\
        Aparelho real & desempenho e hardware verdadeiros & cada modelo exige um aparelho \\
        \bottomrule
      \end{tabularx}
    \end{column}
    \begin{column}{0.48\textwidth}
      \begin{definicao}[Emulador e simulador]
        O \textbf{emulador} imita o aparelho, inclusive o processador e o sistema completo. O \textbf{simulador} apenas reproduz o comportamento do sistema, executando o app compilado para o computador.
      \end{definicao}
      \begin{pordentro}[Por que o emulador exige virtualização]
        O emulador Android roda uma imagem completa do sistema numa máquina virtual. Com aceleração por hardware (Intel VT-x, AMD-V, KVM no Linux, Hypervisor no macOS), as instruções rodam direto no processador e o emulador fica rápido.
      \end{pordentro}
    \end{column}
  \end{columns}
\end{frame}

\begin{frame}{Depuração}
  \begin{columns}[T]
    \begin{column}{0.5\textwidth}
      \begin{definicao}[Depurador]
        Ferramenta que pausa o programa num ponto escolhido (\emph{breakpoint}) para inspecionar variáveis, a pilha de chamadas e executar passo a passo.
      \end{definicao}
      \scriptsize
      \begin{tabularx}{\linewidth}{@{}l L@{}}
        \toprule
        \cabecalho{Comando} & \cabecalho{O que faz} \\
        \midrule
        \emph{Step over} & executa a linha e para na próxima \\
        \emph{Step into} & entra na função chamada \\
        \emph{Step out} & termina a função atual e volta \\
        \emph{Resume} & continua até o próximo \emph{breakpoint} \\
        \bottomrule
      \end{tabularx}
    \end{column}
    \begin{column}{0.48\textwidth}
      \begin{conceito}[Método científico da depuração]
        \begin{enumerate}
          \item Reproduza o problema de forma confiável.
          \item Formule uma hipótese sobre a causa.
          \item Teste a hipótese (\emph{breakpoint}, \emph{log}, teste automatizado).
          \item Corrija e escreva um teste que impeça a volta do defeito.
        \end{enumerate}
      \end{conceito}
      \begin{dica}
        Leia a mensagem de erro e a pilha de chamadas inteiras: a primeira linha do \textbf{seu} código costuma apontar a causa.
      \end{dica}
    \end{column}
  \end{columns}
\end{frame}


% =============================================================================
\section{Flutter e Dart}

\begin{frame}{O que é o Flutter}
  \begin{columns}[c]
    \begin{column}{0.52\textwidth}
      \begin{itemize}
        \item \termo{Flutter}: kit de desenvolvimento de interfaces do Google, de código aberto.
        \item \textbf{Um único código} em \termo{Dart} gera apps para Android, iOS, web, Windows, macOS e Linux.
        \item A interface é feita de \termo{widgets} desenhados pelo próprio motor do Flutter, idênticos em todas as plataformas.
        \item Compilado para código nativo (AOT) em release: desempenho próximo ao nativo.
        \item \emph{Hot reload}: mudanças aparecem no aparelho em menos de um segundo.
      \end{itemize}
    \end{column}
    \begin{column}{0.46\textwidth}
      \centering
      \begin{tikzpicture}[node distance=1.5mm]
        \node[caixaambarforte,minimum width=5cm] (app) {Seu app (Dart + widgets)};
        \node[caixaescura,minimum width=5cm,below=of app] (fw) {Framework Flutter\\ \tiny Material, Cupertino, widgets, animação};
        \node[caixapetroleo,minimum width=5cm,below=of fw] (eng) {Motor (Impeller)\\ \tiny desenha cada pixel, Dart runtime};
        \node[caixa,minimum width=5cm,below=of eng] (emb) {\emph{Embedder} da plataforma\\ \tiny Android, iOS, web, desktop};
        \node[caixa,fill=cinzaclaro,minimum width=5cm,below=of emb] (so) {Sistema operacional e hardware};
      \end{tikzpicture}
    \end{column}
  \end{columns}
  \begin{dica}[Por que Dart?]
    JIT no desenvolvimento (\emph{hot reload}), AOT no release (desempenho), \emph{null safety} e sintaxe familiar a quem conhece Java, C\# ou JavaScript.
  \end{dica}
\end{frame}

% =============================================================================
\section{Instalação}

\begin{frame}[fragile]{Instalando o Flutter SDK}
  \begin{columns}[T]
    \begin{column}{0.5\textwidth}
      \begin{enumerate}
        \item Instale o \textbf{VS Code} e a extensão \textbf{Flutter}; o comando \emph{Flutter: New Project} oferece baixar o SDK.
        \item Ou baixe o SDK no site oficial e ponha \arq{flutter/bin} no \cod{PATH}.
        \item Instale o \textbf{Android Studio} (Android SDK e emulador).
        \item \cod{flutter doctor --android-licenses}.
      \end{enumerate}
      \begin{atencao}[iOS]
        Compilar para iOS exige \textbf{macOS} com Xcode. No Windows e no Linux, teste no Android, na web e no desktop.
      \end{atencao}
    \end{column}
    \begin{column}{0.48\textwidth}
\begin{terminal}[title=flutter doctor (saída resumida)]
$ flutter doctor
[ok] Flutter (Channel stable)
[ok] Android toolchain - Android SDK
[ x] Xcode (somente no macOS)
[ok] Chrome - develop for the web
[ok] Android Studio
[ok] VS Code
[ok] Connected device (2 available)
[ok] Network resources
\end{terminal}
      \begin{dica}
        \cod{flutter doctor -v} mostra o que falta corrigir. Use o canal \textbf{stable}; \cod{flutter upgrade} atualiza o SDK.
      \end{dica}
    \end{column}
  \end{columns}
\end{frame}

\begin{frame}[fragile]{Onde executar o app}
  \begin{columns}[T]
    \begin{column}{0.5\textwidth}
      \scriptsize
      \begin{tabularx}{\linewidth}{@{}l L@{}}
        \toprule
        \cabecalho{Destino} & \cabecalho{Como} \\
        \midrule
        Emulador Android & Device Manager do Android Studio \\
        Celular Android & opções do desenvolvedor + depuração USB ou Wi-Fi \\
        Simulador iOS & Xcode (somente macOS) \\
        iPhone real & Xcode + conta Apple (gratuita para testes) \\
        Navegador & Chrome, sem configuração extra \\
        Desktop & Windows, macOS ou Linux, com as ferramentas de C++ do sistema \\
        \bottomrule
      \end{tabularx}
      \begin{atencao}[Hardware real]
        Bluetooth, NFC, câmera de verdade e sensores (Módulo 13) exigem \textbf{celular físico}.
      \end{atencao}
    \end{column}
    \begin{column}{0.48\textwidth}
\begin{terminal}[title=Dispositivos pelo terminal]
flutter devices              # lista destinos
flutter emulators            # lista AVDs
flutter emulators --launch Pixel_8
flutter run                  # escolhe o destino
flutter run -d chrome        # roda na web
flutter run -d emulator-5554
\end{terminal}
      \begin{dica}[Celular Android]
        Ative as \textbf{Opções do desenvolvedor} (toque 7 vezes em \emph{Número da versão}) e a \textbf{Depuração USB}. O mesmo vale para a trilha Kotlin (Módulo 01).
      \end{dica}
    \end{column}
  \end{columns}
\end{frame}

% =============================================================================
\section{Estrutura do projeto}

\begin{frame}[fragile]{Criando o projeto}
  \begin{columns}[T]
    \begin{column}{0.5\textwidth}
\begin{terminal}[title=Criando pelo terminal]
flutter create --org br.edu.ifma \
  --platforms android,ios,web labmonitor
cd labmonitor
flutter run
\end{terminal}
\begin{codigo}[title=labmonitor/,listing options={style=dm,language=bash,basicstyle=\ttfamily\tiny}]{bash}
labmonitor/
|-- lib/
|   `-- main.dart          ponto de entrada
|-- test/
|   `-- widget_test.dart   testes
|-- android/               projeto Android (Gradle)
|-- ios/                   projeto iOS (Xcode)
|-- web/
|-- pubspec.yaml           dependências e recursos
|-- pubspec.lock           versões exatas
`-- analysis_options.yaml  regras de análise
\end{codigo}
    \end{column}
    \begin{column}{0.48\textwidth}
      \begin{conceito}[\cod{--org}]
        Define o identificador do app: \cod{br.edu.ifma.labmonitor}. Ele vira o \cod{applicationId} no Android e o \emph{bundle identifier} no iOS, e não pode mudar depois de publicado.
      \end{conceito}
      \begin{itemize}
        \item Quase todo o seu código fica em \arq{lib/}.
        \item \arq{android/} e \arq{ios/} são projetos nativos completos: você mexe neles para permissões, ícones e assinatura.
        \item \arq{pubspec.lock} vai para o Git; garante as mesmas versões para toda a equipe.
      \end{itemize}
    \end{column}
  \end{columns}
\end{frame}

\begin{frame}[fragile]{\arq{pubspec.yaml}}
  \begin{columns}[T]
    \begin{column}{0.5\textwidth}
\begin{codigo}[title=pubspec.yaml]{YAML}
name: labmonitor
version: 1.0.0+1          # nome+número do build

environment:
  sdk: ^3.9.0

dependencies:
  flutter:
    sdk: flutter
  go_router: ^16.0.0      # exemplo de pacote

dev_dependencies:
  flutter_test:
    sdk: flutter
  flutter_lints: ^6.0.0

flutter:
  assets: [assets/imagens/]
\end{codigo}
    \end{column}
    \begin{column}{0.48\textwidth}
      \begin{itemize}
        \item \cod{version}: \cod{1.0.0} é o nome exibido; \cod{+1} é o número do build (\cod{versionCode} no Android).
        \item \cod{dependencies}: pacotes usados pelo app.
        \item \cod{dev\_dependencies}: só para desenvolvimento e testes.
        \item \cod{assets}: imagens e arquivos empacotados (as versões do exemplo são ilustrativas).
      \end{itemize}
      \begin{atencao}[Indentação]
        YAML usa espaços para definir a hierarquia. Um espaço a mais ou a menos quebra o arquivo.
      \end{atencao}
    \end{column}
  \end{columns}
\end{frame}

\begin{frame}[fragile]{\arq{lib/main.dart}: o primeiro app}
  \begin{columns}[T]
    \begin{column}{0.56\textwidth}
\begin{codigo}[title=lib/main.dart]{Dart}
import 'package:flutter/material.dart';

void main() {
  runApp(const LabMonitorApp());
}

class LabMonitorApp extends StatelessWidget {
  const LabMonitorApp({super.key});
  @override
  Widget build(BuildContext context) {
    return MaterialApp(
      title: 'LabMonitor',
      theme: ThemeData(
        colorSchemeSeed: const Color(0xFF0D3052)),
      home: Scaffold(
        appBar: AppBar(title: const Text('LabMonitor')),
        body: const Center(child: Text('Olá, laboratório!')),
      ),
    );
  }
}
\end{codigo}
    \end{column}
    \begin{column}{0.42\textwidth}
      \begin{itemize}
        \item \cod{main()} chama \cod{runApp} com o widget raiz.
        \item \cod{StatelessWidget}: widget sem estado próprio; descreve a tela no método \cod{build}.
        \item \cod{MaterialApp}: tema, rotas e localização do app.
        \item \cod{Scaffold}: estrutura de tela com barra superior e corpo.
        \item \cod{const}: widget criado em tempo de compilação, sem custo de reconstrução.
      \end{itemize}
    \end{column}
  \end{columns}
\end{frame}

% =============================================================================
\section{Executando e depurando}

\begin{frame}{\emph{Hot reload} e modos de execução}
  \begin{columns}[T]
    \begin{column}{0.5\textwidth}
      \scriptsize
      \begin{tabularx}{\linewidth}{@{}l L L@{}}
        \toprule
        \cabecalho{Ação} & \cabecalho{Tecla} & \cabecalho{Efeito} \\
        \midrule
        \emph{Hot reload} & \cod{r} & injeta o código novo e mantém o estado \\
        \emph{Hot restart} & \cod{R} & reinicia o app, perde o estado \\
        \emph{Full restart} & parar e rodar & necessário após mudar código nativo ou o \arq{pubspec} \\
        \bottomrule
      \end{tabularx}
      \begin{dica}
        No VS Code, salvar o arquivo já dispara o \emph{hot reload}.
      \end{dica}
    \end{column}
    \begin{column}{0.48\textwidth}
      \scriptsize
      \begin{tabularx}{\linewidth}{@{}l L@{}}
        \toprule
        \cabecalho{Modo} & \cabecalho{Uso} \\
        \midrule
        \textbf{debug} & desenvolvimento: JIT, \emph{hot reload}, asserções, mais lento \\
        \textbf{profile} & medir desempenho num celular real (\cod{flutter run --profile}) \\
        \textbf{release} & publicação: AOT, otimizado, sem ferramentas de depuração \\
        \bottomrule
      \end{tabularx}
      \begin{atencao}
        Nunca avalie desempenho em modo debug: ele pode ser várias vezes mais lento que o release.
      \end{atencao}
    \end{column}
  \end{columns}
\end{frame}

\begin{frame}[fragile]{Flutter DevTools e logs}
  \begin{columns}[T]
    \begin{column}{0.5\textwidth}
      \begin{conceito}[DevTools]
        Ferramentas no navegador ou no editor:
        \begin{itemize}
          \item \textbf{Inspector}: árvore de widgets e depurador de layout.
          \item \textbf{Performance}: quadros lentos (\emph{jank}).
          \item \textbf{CPU} e \textbf{Memory}: perfis e vazamentos.
          \item \textbf{Network}: requisições HTTP do app.
          \item \textbf{Logging}: mensagens e erros.
        \end{itemize}
      \end{conceito}
    \end{column}
    \begin{column}{0.48\textwidth}
\begin{codigo}[title=Registrando logs]{Dart}
import 'dart:developer' as dev;

void lerSensor(double valor) {
  debugPrint('Leitura: $valor');
  if (valor > 80) {
    dev.log('Temperatura alta: $valor',
        name: 'LabMonitor', level: 900);
  }
}
\end{codigo}
      \begin{dica}[Depurador]
        No VS Code: \emph{breakpoint} na margem e \textbf{F5} para depurar.
      \end{dica}
      \begin{atencao}
        Nunca registre senhas, tokens ou dados pessoais.
      \end{atencao}
    \end{column}
  \end{columns}
\end{frame}

% =============================================================================
\section{Pacotes e comandos}

\begin{frame}[fragile]{Pacotes do pub.dev}
  \begin{columns}[T]
    \begin{column}{0.5\textwidth}
\begin{terminal}[title=Gerenciando dependências]
flutter pub add go_router     # adiciona
flutter pub add --dev mocktail
flutter pub remove go_router
flutter pub get               # baixa tudo
flutter pub outdated          # o que tem versão nova
flutter pub upgrade
\end{terminal}
      \begin{conceito}[Versões]
        \cod{\^{}2.1.0} aceita qualquer \cod{2.x.y} a partir de \cod{2.1.0}, mas não \cod{3.0.0}, que pode quebrar a compatibilidade (versionamento semântico).
      \end{conceito}
    \end{column}
    \begin{column}{0.48\textwidth}
      \begin{dica}[Como escolher um pacote]
        \begin{itemize}
          \item \textbf{Publisher verificado} (ex.: \cod{flutter.dev}, \cod{firebase.google.com}).
          \item \textbf{Pub points} altos e atualizações recentes.
          \item Plataformas suportadas (Android, iOS, web\ldots).
          \item \textbf{Licença} compatível com o seu uso.
        \end{itemize}
      \end{dica}
      \begin{atencao}
        Pacote é código de terceiros no seu app: prefira poucos e bem mantidos.
      \end{atencao}
    \end{column}
  \end{columns}
\end{frame}

\begin{frame}{Comandos do dia a dia}
  \footnotesize
  \begin{tabularx}{\textwidth}{@{}l L@{}}
    \toprule
    \cabecalho{Comando} & \cabecalho{Para que serve} \\
    \midrule
    \cod{flutter doctor} & diagnostica o ambiente \\
    \cod{flutter create} & cria um projeto \\
    \cod{flutter run} & compila, instala e executa com \emph{hot reload} \\
    \cod{flutter analyze} & análise estática (erros e más práticas) \\
    \cod{dart format .} & formata todo o código \\
    \cod{flutter test} & roda os testes (Módulo 14) \\
    \cod{flutter build apk} / \cod{appbundle} / \cod{ipa} / \cod{web} & gera os pacotes de publicação (Módulo 15) \\
    \cod{flutter clean} & apaga a pasta \arq{build/} quando algo ``estranho'' acontecer \\
    \cod{dart run build\_runner build} & executa geradores de código (JSON, Drift, Riverpod) \\
    \bottomrule
  \end{tabularx}
  \begin{dica}[Git]
    O \arq{.gitignore} gerado pelo \cod{flutter create} já exclui \arq{build/} e \arq{.dart\_tool/}. Faça o primeiro commit e publique no GitHub (Git em detalhes no Módulo 03).
  \end{dica}
\end{frame}

% =============================================================================
\section{Fechamento}

\begin{frame}{Resumo}
  \begin{resumo}
    \begin{itemize}
      \item Flutter desenha a interface com o próprio motor; um código Dart atende Android, iOS, web e desktop.
      \item \cod{flutter doctor} confere o ambiente; iOS exige macOS e Xcode.
      \item O código vive em \arq{lib/}; \arq{pubspec.yaml} declara versão, dependências e recursos.
      \item \emph{Hot reload} (\cod{r}) mantém o estado; \emph{hot restart} (\cod{R}) reinicia.
      \item Meça desempenho em modo \textbf{profile}; publique em \textbf{release}.
      \item Escolha pacotes do pub.dev por publisher, manutenção, plataformas e licença.
    \end{itemize}
  \end{resumo}
\end{frame}

\begin{frame}{Glossário}
  \begin{glossario}
    SDK & conjunto de bibliotecas e ferramentas para desenvolver para uma plataforma \\
    Flutter SDK & \emph{framework}, motor, Dart SDK e a ferramenta \cod{flutter} \\
    \cod{flutter doctor} & diagnóstico do ambiente: SDKs, licenças, aparelhos \\
    \cod{pubspec.yaml} & arquivo com nome, versão, dependências e recursos do projeto \\
    pub.dev & repositório oficial de pacotes Dart e Flutter \\
    \emph{Hot reload} & injeta o código alterado no app em execução mantendo o estado \\
    \emph{Hot restart} & reinicia o app com o código novo, perdendo o estado \\
    DevTools & ferramentas de inspeção, desempenho, memória e rede do Flutter \\
    \emph{Breakpoint} & ponto em que o depurador pausa a execução \\
  \end{glossario}
\end{frame}

\begin{frame}{Questões de revisão}
  \begin{questoes}
    \item Qual a diferença entre JIT e AOT? Como o Flutter usa cada um?
    \item Por que o \emph{hot reload} é possível em modo \emph{debug} e não em \emph{release}?
    \item O que o \cod{flutter doctor} verifica e por que rodá-lo depois de instalar?
    \item Para que serve o \cod{pubspec.lock} e por que ele vai para o Git?
    \item Explique a diferença entre emulador Android e simulador iOS.
    \item Cite duas situações em que é preciso um aparelho real.
    \item Quando usar \emph{hot restart} em vez de \emph{hot reload}?
    \item Descreva o método para depurar um app que fecha ao tocar num botão.
  \end{questoes}
\end{frame}

\begin{frame}{Exercícios}
  \begin{exercicio}[Exercício 1 --- Ambiente]
    Rode \cod{flutter doctor -v}, corrija todos os itens possíveis no seu sistema e execute o app padrão no emulador, no celular e no Chrome.
  \end{exercicio}
  \begin{exercicio}[Exercício 2 --- Hot reload]
    Mude a cor-semente do tema e o texto da tela com o app rodando. Compare \emph{hot reload} e \emph{hot restart} usando o contador do app padrão.
  \end{exercicio}
  \begin{exercicio}[Exercício 3 --- Pacotes]
    Pesquise no pub.dev três pacotes para gráficos. Compare publisher, pub points, plataformas e licença e justifique qual usaria.
  \end{exercicio}
  \begin{projeto}[LabMonitor --- Etapa 1]
    Crie o projeto \cod{labmonitor} (\cod{--org br.edu.ifma}), troque o texto por ``LabMonitor --- DEE/IFMA'', adicione um \cod{debugPrint} no \cod{main} e publique o repositório no GitHub.
  \end{projeto}
\end{frame}

\begin{frame}{Para aprofundar}
  \small
  \begin{itemize}
    \item Instalação: \url{https://docs.flutter.dev/get-started}
    \item Arquitetura do Flutter: \url{https://docs.flutter.dev/resources/architectural-overview}
    \item Hot reload: \url{https://docs.flutter.dev/tools/hot-reload}
    \item DevTools: \url{https://docs.flutter.dev/tools/devtools}
    \item Usando pacotes: \url{https://docs.flutter.dev/packages-and-plugins/using-packages}
    \item Repositório de pacotes: \url{https://pub.dev}
  \end{itemize}
\end{frame}

\begin{frame}{Respostas das questões de revisão}
  \begin{respostas}
    \item JIT compila durante a execução; AOT, antes. O Flutter usa JIT em \emph{debug} (para o \emph{hot reload}) e AOT em \emph{profile} e \emph{release} (código ARM rápido).
    \item O \emph{hot reload} depende da VM do Dart com JIT, que aceita código novo em execução; em \emph{release} o código já foi compilado AOT.
    \item Instalação do Flutter, Android SDK e licenças, Xcode (no macOS), IDE e aparelhos conectados; aponta o que falta configurar.
    \item Registra as versões exatas resolvidas de todas as dependências, para que todos compilem o mesmo app.
    \item O emulador Android imita o aparelho e roda o sistema completo numa máquina virtual; o simulador iOS executa o app compilado para o Mac, sem imitar o hardware.
    \item Bluetooth, NFC e USB; medir desempenho e bateria; câmera e sensores reais.
    \item Quando a mudança afeta o estado inicial (\cod{main}, \cod{initState}, variáveis globais) ou a estrutura de classes, que o \emph{hot reload} não reaplica.
    \item Reproduzir o fechamento, ler a pilha no console, colocar um \emph{breakpoint} antes da linha suspeita, testar a hipótese, corrigir e criar um teste.
  \end{respostas}
\end{frame}

\encerramento{Módulo 02 --- Dart essencial: tipos, null safety, funções, classes, coleções, records e padrões.}

\end{document}
